% ============================================================
%  RAPPORT TECHNIQUE N°3 - Deep Learning
%  Backpropagation dans la Couche de Convolution
%  Auteurs : Souaibou DIENG & Ibrahima FALL
%  ESP DAKAR - Master 2 SRT - 2025/2026
% ============================================================

\documentclass[12pt, a4paper]{article}

% ---------- Encodage & Langue ----------
\usepackage[utf8]{inputenc}
\usepackage[T1]{fontenc}
\usepackage[french]{babel}

% ---------- Mise en page ----------
\usepackage[top=2.5cm, bottom=2.5cm, left=3cm, right=2.5cm]{geometry}
\usepackage{setspace}
\onehalfspacing

% ---------- Polices ----------
\usepackage{lmodern}
\usepackage{microtype}

% ---------- Graphiques & Couleurs ----------
\usepackage{graphicx}
\usepackage{xcolor}
\usepackage{float}
\usepackage{tikz}
\usetikzlibrary{matrix, arrows.meta, positioning, fit, backgrounds}

% Palette de couleurs
\definecolor{espblue}{RGB}{0, 70, 127}
\definecolor{espgold}{RGB}{204, 153, 0}
\definecolor{codegray}{RGB}{40, 44, 52}
\definecolor{codegreen}{RGB}{80, 200, 120}
\definecolor{codeorange}{RGB}{230, 126, 34}
\definecolor{codecomment}{RGB}{127, 140, 141}
\definecolor{codered}{RGB}{231, 76, 60}
\definecolor{cellblue}{RGB}{219, 234, 254}
\definecolor{cellgold}{RGB}{254, 243, 199}
\definecolor{cellgreen}{RGB}{209, 250, 229}

% ---------- Tableaux ----------
\usepackage{booktabs}
\usepackage{tabularx}
\usepackage{array}
\usepackage{colortbl}
\usepackage{nicematrix}

% ---------- Maths ----------
\usepackage{amsmath}
\usepackage{amssymb}
\usepackage{amsfonts}

% ---------- Code source ----------
\usepackage{listings}
\lstdefinestyle{pythonstyle}{
    backgroundcolor=\color{codegray},
    basicstyle=\ttfamily\footnotesize\color{white},
    keywordstyle=\color{codeorange}\bfseries,
    stringstyle=\color{codegreen},
    commentstyle=\color{codecomment}\itshape,
    numberstyle=\tiny\color{gray},
    numbers=left,
    numbersep=8pt,
    breaklines=true,
    breakatwhitespace=true,
    frame=single,
    framesep=5pt,
    rulecolor=\color{espblue},
    showstringspaces=false,
    tabsize=4,
    captionpos=b,
    literate=
        {é}{{\'{e}}}1 {è}{{\`{e}}}1 {ê}{{\^{e}}}1 {ë}{{\"e}}1
        {à}{{\`{a}}}1 {â}{{\^{a}}}1 {î}{{\^i}}1 {ô}{{\^o}}1
        {ù}{{\`{u}}}1 {û}{{\^{u}}}1 {ç}{{\c{c}}}1
}
\lstset{style=pythonstyle}

% ---------- Hyperliens ----------
\usepackage[hidelinks]{hyperref}
\hypersetup{
    pdftitle={Rapport TP3 - Backpropagation dans la couche de convolution},
    pdfauthor={Souaibou DIENG \& Ibrahima FALL},
}

% ---------- En-têtes et pieds de page ----------
\usepackage{fancyhdr}
\pagestyle{fancy}
\fancyhf{}
\fancyhead[L]{\small\textcolor{espblue}{\textit{Backpropagation dans la Couche de Convolution}}}
\fancyhead[R]{\small\textcolor{espblue}{\textit{Master 2 SRT -- ESP Dakar}}}
\fancyfoot[C]{\textcolor{espblue}{\thepage}}
\renewcommand{\headrulewidth}{0.4pt}
\renewcommand{\headrule}{\hbox to\headwidth{%
    \color{espblue}\leaders\hrule height \headrulewidth\hfill}}

% ---------- Boîtes ----------
\usepackage{tcolorbox}
\tcbuselibrary{skins, breakable}

\newtcolorbox{infobox}[1][]{
    enhanced, breakable,
    colback=espblue!5, colframe=espblue,
    boxrule=0.8pt, arc=4pt,
    fonttitle=\bfseries\color{white},
    colbacktitle=espblue, #1
}

\newtcolorbox{mathbox}{
    enhanced,
    colback=espgold!8, colframe=espgold,
    boxrule=0.8pt, arc=4pt,
}

\newtcolorbox{resultbox}{
    enhanced,
    colback=codegreen!8, colframe=codegreen!60!black,
    boxrule=0.8pt, arc=4pt,
}

\newtcolorbox{stepbox}[1]{
    enhanced, breakable,
    colback=espblue!4, colframe=espblue!35,
    boxrule=0.6pt, arc=3pt,
    title={\textbf{#1}},
    fonttitle=\small\bfseries\color{espblue},
    colbacktitle=espblue!12,
    coltitle=espblue,
}

\newtcolorbox{propbox}{
    enhanced,
    colback=codered!5, colframe=codered!50,
    boxrule=0.6pt, arc=3pt,
}

% ---------- Sections personnalisées ----------
\usepackage{titlesec}
\titleformat{\section}
    {\normalfont\Large\bfseries\color{espblue}}
    {\thesection.}{10pt}{}
    [\vspace{2pt}{\color{espblue}\titlerule[0.6pt]}]
\titlespacing{\section}{0pt}{18pt}{10pt}

\titleformat{\subsection}
    {\normalfont\large\bfseries\color{espblue!80!black}}
    {\thesubsection}{8pt}{}
\titlespacing{\subsection}{0pt}{12pt}{6pt}

% ---------- Divers ----------
\usepackage{enumitem}
\usepackage{caption}
\captionsetup{font=small, labelfont=bf, labelsep=period}

% ============================================================
%  DÉBUT DU DOCUMENT
% ============================================================
\begin{document}

% ============================================================
%  PAGE DE GARDE
% ============================================================
\begin{titlepage}
    \centering
    \vspace*{0.3cm}

    {\color{espblue}\rule{\linewidth}{3pt}}\\[0.5cm]

    \begin{minipage}[c]{0.30\textwidth}
        \centering
        {\small\bfseries RÉPUBLIQUE DU SÉNÉGAL}\\[3pt]
        {\footnotesize Un Peuple – Un But – Une Foi}
    \end{minipage}
    \hfill
    \begin{minipage}[c]{0.30\textwidth}
        \centering
        \includegraphics[width=2.8cm]{logo_esp.png}
    \end{minipage}
    \hfill
    \begin{minipage}[c]{0.30\textwidth}
        \centering
        {\small\bfseries UNIVERSITÉ CHEIKH ANTA DIOP}\\[3pt]
        {\footnotesize DE DAKAR (UCAD)}
    \end{minipage}

    \vspace{0.5cm}
    {\color{espblue}\rule{\linewidth}{1pt}}\\[0.3cm]

    {\large\bfseries\color{espblue} ÉCOLE SUPÉRIEURE POLYTECHNIQUE DE DAKAR}\\[4pt]
    {\normalsize DÉPARTEMENT DE GÉNIE INFORMATIQUE}\\[2pt]
    {\normalsize\bfseries MASTER 2 -- SYSTÈMES, RÉSEAUX ET TÉLÉCOMMUNICATIONS}

    \vspace{0.4cm}
    {\color{espgold}\rule{0.5\linewidth}{2pt}}
    \vspace{0.3cm}

    {\large\color{espblue!70!black} Module :}\\[3pt]
    {\Large\bfseries Deep Learning / Traitement d'Images}

    \vspace{0.3cm}
    {\color{espgold}\rule{0.8\linewidth}{1.5pt}}
    \vspace{0.3cm}

    % Badge numéro de TP
    \begin{tcolorbox}[enhanced, width=4cm, colback=espblue, colframe=espblue,
        arc=20pt, boxrule=0pt, halign=center]
        {\large\bfseries\color{white} Rapport Technique}\\[2pt]
        {\LARGE\bfseries\color{espgold} N° 3}
    \end{tcolorbox}

    \vspace{0.3cm}

    {\LARGE\bfseries\color{espblue}
        Backpropagation dans la\\[5pt]
        Couche de Convolution
    }

    \vspace{0.3cm}
    {\color{espgold}\rule{0.8\linewidth}{1.5pt}}
    \vspace{0.5cm}

    \Large\bfseries Réalisé par :\\[0.4cm]
    \begin{tabular}{cc}
        \Large\color{espblue} Souaibou \textsc{Dieng} &
        \Large\color{espblue} Ibrahima \textsc{Fall} \\
    \end{tabular}

    \vfill

    \begin{tabular}{rl}
        \textbf{Classe :}              & Master 2 SRT \\[3pt]
        \textbf{Encadrant :}           & Dr. Mouhamed Kaba \textsc{Keïta} \\[3pt]
        \textbf{Année universitaire :} & 2025 -- 2026 \\
    \end{tabular}

    \vspace{0.4cm}
    {\color{espblue}\rule{\linewidth}{3pt}}
\end{titlepage}

% ============================================================
%  TABLE DES MATIÈRES
% ============================================================
\tableofcontents
\thispagestyle{fancy}
\newpage

% ============================================================
%  INTRODUCTION
% ============================================================
\section*{Introduction}
\addcontentsline{toc}{section}{Introduction}

L'efficacité des \textbf{réseaux de neurones convolutifs} (CNN) repose sur le
\textbf{partage de poids}, ce qui leur confère une invariance spatiale essentielle pour
le traitement des images. Cependant, leur entraînement nécessite une adaptation
spécifique de l'algorithme de rétropropagation.

Ce rapport analyse \textbf{mathématiquement} comment les gradients de la fonction
de perte $L$ circulent à travers une opération de convolution entre une entrée $X$
et un filtre $F$, en distinguant :
\begin{enumerate}[leftmargin=2cm, label=\textbf{\arabic{*}.}]
    \item La propagation avant (\textit{Forward Pass}) et les équations de convolution.
    \item Les gradients locaux par rapport au filtre et à l'entrée.
    \item La rétropropagation complète de la perte $L$.
    \item Une implémentation Python de validation.
\end{enumerate}

% ============================================================
%  SECTION 1 : FORWARD PASS
% ============================================================
\section{Propagation Avant (\textit{Forward Pass})}

\subsection{Dimensionnement}

On considère une convolution avec un \textbf{pas (\textit{stride}) de 1} et sans
\textbf{remplissage (\textit{padding})}. Les dimensions sont :

\begin{table}[H]
    \centering
    \caption{Dimensions des tenseurs de la convolution}
    \label{tab:dims}
    \renewcommand{\arraystretch}{1.4}
    \begin{tabular}{lccc}
        \toprule
        \rowcolor{espblue!15}
        \textbf{Tenseur} & \textbf{Notation} & \textbf{Taille} & \textbf{Rôle} \\
        \midrule
        Entrée   & $X$ & $3 \times 3$ & Image d'entrée \\
        Filtre   & $F$ & $2 \times 2$ & Noyau de convolution \\
        Sortie   & $O$ & $2 \times 2$ & Carte de caractéristiques \\
        \bottomrule
    \end{tabular}
\end{table}

\subsection{Visualisation de l'Opération}

\begin{figure}[H]
\centering
\begin{tikzpicture}[
    cell/.style={draw=espblue!50, minimum size=0.85cm, font=\small},
    highlight/.style={draw=espblue, fill=cellblue, minimum size=0.85cm, font=\small\bfseries},
    filter/.style={draw=espgold, fill=cellgold, minimum size=0.85cm, font=\small\bfseries},
    out/.style={draw=codegreen!60!black, fill=cellgreen, minimum size=0.85cm, font=\small\bfseries},
]
    % --- Matrice X ---
    \node[font=\small\bfseries\color{espblue}] at (1.3, 3.2) {Entrée $X$ ($3\times3$)};
    \node[highlight] at (0.0, 2.55) {$X_{11}$};
    \node[highlight] at (0.85,2.55) {$X_{12}$};
    \node[cell]      at (1.70,2.55) {$X_{13}$};
    \node[highlight] at (0.0, 1.70) {$X_{21}$};
    \node[highlight] at (0.85,1.70) {$X_{22}$};
    \node[cell]      at (1.70,1.70) {$X_{23}$};
    \node[cell]      at (0.0, 0.85) {$X_{31}$};
    \node[cell]      at (0.85,0.85) {$X_{32}$};
    \node[cell]      at (1.70,0.85) {$X_{33}$};

    % --- Signe * ---
    \node[font=\LARGE\bfseries\color{espblue}] at (2.9, 1.70) {$*$};

    % --- Filtre F ---
    \node[font=\small\bfseries\color{espgold}] at (4.25, 3.2) {Filtre $F$ ($2\times2$)};
    \node[filter] at (3.60, 2.55) {$F_{11}$};
    \node[filter] at (4.45, 2.55) {$F_{12}$};
    \node[filter] at (3.60, 1.70) {$F_{21}$};
    \node[filter] at (4.45, 1.70) {$F_{22}$};

    % --- Signe = ---
    \node[font=\LARGE\bfseries\color{espblue}] at (5.65, 1.70) {$=$};

    % --- Sortie O ---
    \node[font=\small\bfseries\color{codegreen!60!black}] at (7.15, 3.2) {Sortie $O$ ($2\times2$)};
    \node[out] at (6.60, 2.55) {$O_{11}$};
    \node[out] at (7.45, 2.55) {$O_{12}$};
    \node[out] at (6.60, 1.70) {$O_{21}$};
    \node[out] at (7.45, 1.70) {$O_{22}$};
\end{tikzpicture}
\caption{Schéma de la convolution $3\times3 * 2\times2 \to 2\times2$. Les cellules bleues
indiquent la position du filtre pour le calcul de $O_{11}$.}
\label{fig:conv}
\end{figure}

\subsection{Équations de la Convolution}

La formule générale est $O_{ij} = \sum_{p,q} X_{i+p,\,j+q} \cdot F_{pq}$.
Développée pour chaque sortie :

\begin{mathbox}
\begin{align}
    O_{11} &= X_{11}F_{11} + X_{12}F_{12} + X_{21}F_{21} + X_{22}F_{22} \label{eq:o11}\\
    O_{12} &= X_{12}F_{11} + X_{13}F_{12} + X_{22}F_{21} + X_{23}F_{22} \label{eq:o12}\\
    O_{21} &= X_{21}F_{11} + X_{22}F_{12} + X_{31}F_{21} + X_{32}F_{22} \label{eq:o21}\\
    O_{22} &= X_{22}F_{11} + X_{23}F_{12} + X_{32}F_{21} + X_{33}F_{22} \label{eq:o22}
\end{align}
\end{mathbox}

\begin{infobox}[title={Observation clé : le partage de poids}]
Le \textbf{même filtre} $F$ est appliqué à chaque position. Chaque élément $F_{pq}$
contribue aux \textbf{quatre sorties} $O_{11}, O_{12}, O_{21}, O_{22}$.
C'est ce partage de poids qui donne aux CNN leur efficacité et leur invariance spatiale.
\end{infobox}

% ============================================================
%  SECTION 2 : GRADIENTS LOCAUX
% ============================================================
\section{Calcul des Gradients Locaux}

Le gradient local représente la sensibilité d'une sortie à une variation
d'un paramètre donné (\textit{à la}  perte non encore intégrée).

\subsection{Dérivées par rapport au Filtre $F$}

En différenciant les équations~\eqref{eq:o11}--\eqref{eq:o22}, pour la sortie
$O_{12}$ :

\begin{stepbox}{Gradients locaux de $O_{12}$ par rapport à $F$}
\begin{equation}
    \frac{\partial O_{12}}{\partial F_{11}} = X_{12},\quad
    \frac{\partial O_{12}}{\partial F_{12}} = X_{13},\quad
    \frac{\partial O_{12}}{\partial F_{21}} = X_{22},\quad
    \frac{\partial O_{12}}{\partial F_{22}} = X_{23}
\end{equation}
\textit{Interprétation :} le gradient local par rapport à $F_{pq}$ est simplement
l'élément d'entrée $X$ qui était « sous » $F_{pq}$ lors du calcul de $O_{12}$.
\end{stepbox}

\subsection{Dérivées par rapport à l'Entrée $X$}

Certaines entrées n'interviennent pas dans le calcul de sorties spécifiques.
Pour $O_{12}$ :

\begin{stepbox}{Gradients locaux de $O_{12}$ par rapport à $X$}
\begin{equation}
    \frac{\partial O_{12}}{\partial X_{11}} = 0,\quad
    \frac{\partial O_{12}}{\partial X_{12}} = F_{11},\quad
    \frac{\partial O_{12}}{\partial X_{22}} = F_{21}
\end{equation}
\textit{Interprétation :} $X_{11}$ n'apparaît pas dans $O_{12}$, d'où un gradient nul.
\end{stepbox}

% ============================================================
%  SECTION 3 : RÉTROPROPAGATION
% ============================================================
\section{Gradients de la Fonction de Perte $L$}

En appliquant la \textbf{règle de la chaîne}, on accumule les gradients sur
toutes les sorties influencées par un paramètre donné.

\subsection{Gradient par rapport au Filtre $\partial L / \partial F$}

L'élément $F_{12}$ contribue aux quatre sorties. Son gradient total est :

\begin{stepbox}{Gradient total de $L$ par rapport à $F_{12}$}
\begin{equation}
    \frac{\partial L}{\partial F_{12}}
    = \frac{\partial L}{\partial O_{11}} \cdot X_{12}
    + \frac{\partial L}{\partial O_{12}} \cdot X_{13}
    + \frac{\partial L}{\partial O_{21}} \cdot X_{22}
    + \frac{\partial L}{\partial O_{22}} \cdot X_{23}
    \label{eq:grad_F12}
\end{equation}
\end{stepbox}

\begin{propbox}
\textbf{Propriété générale :}
\[
    \frac{\partial L}{\partial F} = \text{Convolution}\!\left(X,\;\frac{\partial L}{\partial O}\right)
\]
Le gradient par rapport au filtre est une \textbf{convolution} de l'entrée $X$
par le gradient de sortie $\delta O = \partial L / \partial O$.
\end{propbox}

\subsection{Gradient par rapport à l'Entrée $\partial L / \partial X$}

Le gradient dépend du nombre de fois que chaque élément de $X$ est « balayé »
par le filtre.

\begin{stepbox}{Cas $X_{11}$ : 1 seule contribution}
\begin{equation}
    \frac{\partial L}{\partial X_{11}}
    = \frac{\partial L}{\partial O_{11}} \cdot F_{11}
\end{equation}
$X_{11}$ n'apparaît que dans $O_{11}$, donc une seule contribution.
\end{stepbox}

\vspace{0.4cm}

\begin{stepbox}{Cas $X_{22}$ : 4 contributions (position centrale)}
\begin{equation}
    \frac{\partial L}{\partial X_{22}}
    = \frac{\partial L}{\partial O_{11}} F_{22}
    + \frac{\partial L}{\partial O_{12}} F_{21}
    + \frac{\partial L}{\partial O_{21}} F_{12}
    + \frac{\partial L}{\partial O_{22}} F_{11}
    \label{eq:grad_x22}
\end{equation}
$X_{22}$ est présent dans les 4 positions de balayage du filtre.
\end{stepbox}

\begin{propbox}
\textbf{Propriété générale :}
\[
    \frac{\partial L}{\partial X}
    = \text{Convolution\_Pleine}\!\left(\frac{\partial L}{\partial O},\;F^{\text{rot180°}}\right)
\]
Le gradient par rapport à $X$ s'obtient par une \textbf{corrélation croisée}
entre $\delta O$ et le filtre \textbf{retourné de 180°}.
\end{propbox}

\begin{table}[H]
    \centering
    \caption{Récapitulatif : nombre de contributions selon la position dans $X$}
    \label{tab:contrib}
    \renewcommand{\arraystretch}{1.5}
    \begin{tabular}{ccc}
        \toprule
        \rowcolor{espblue!15}
        \textbf{Élément de $X$} & \textbf{Contributions} & \textbf{Position} \\
        \midrule
        $X_{11}, X_{13}, X_{31}, X_{33}$ & 1 & Coins \\
        $X_{12}, X_{21}, X_{23}, X_{32}$ & 2 & Bords \\
        $X_{22}$                          & \cellcolor{codegreen!20}\textbf{4} & Centre \\
        \bottomrule
    \end{tabular}
\end{table}

% ============================================================
%  SECTION 4 : IMPLÉMENTATION PYTHON
% ============================================================
\section{Implémentation Logicielle}

Le script suivant simule le calcul du gradient $\partial L / \partial F$
pour une convolution $3 \times 3 \,*\, 2 \times 2$, puis valide le gradient
$\partial L / \partial X$ par corrélation avec $F^{\text{rot180°}}$.

\begin{lstlisting}[language=Python, caption={Calcul des gradients convolutifs -- NumPy}, label=lst:conv]
import numpy as np

# --- Tenseurs initiaux ---
X     = np.array([[1, 2, 3],
                  [4, 5, 6],
                  [7, 8, 9]], dtype=float)

F     = np.array([[1, 0],
                  [0, 1]], dtype=float)

dL_dO = np.ones((2, 2))  # Gradient hypothetique (uniforme)

# =========================================================
# 1. GRADIENT PAR RAPPORT AU FILTRE : dL/dF
# =========================================================
dL_dF = np.zeros_like(F)

for p in range(2):
    for q in range(2):
        # Chaque element du filtre recoit la somme des contributions
        dL_dF[p, q] = np.sum(X[p:p+2, q:q+2] * dL_dO)

print("=== Gradient dL/dF (convolution) ===")
print(dL_dF)

# =========================================================
# 2. GRADIENT PAR RAPPORT A L'ENTREE : dL/dX
# =========================================================
F_rot180 = np.rot90(F, 2)  # Rotation 180 degres du filtre

# Padding complet de dL_dO pour la convolution "pleine"
dL_dO_pad = np.pad(dL_dO, pad_width=1, mode='constant')
dL_dX     = np.zeros_like(X)

for i in range(3):
    for j in range(3):
        dL_dX[i, j] = np.sum(dL_dO_pad[i:i+2, j:j+2] * F_rot180)

print("\n=== Gradient dL/dX (conv. pleine + filtre rot180) ===")
print(dL_dX)
\end{lstlisting}

\begin{resultbox}
\textbf{Résultat attendu pour \texttt{dL/dF}} (avec $\delta O = \mathbf{1}_{2\times2}$) :
les valeurs correspondent à la somme des sous-matrices $2\times2$ de $X$
(position coins, bords, centre), confirmant l'équation~\eqref{eq:grad_F12}.
\end{resultbox}

% ============================================================
%  CONCLUSION
% ============================================================
\section*{Conclusion}
\addcontentsline{toc}{section}{Conclusion}

L'étude de la rétropropagation dans une couche de convolution révèle une
\textbf{symétrie mathématique remarquable} :

\begin{resultbox}
\begin{itemize}[noitemsep, leftmargin=1.5cm]
    \item \textbf{Gradient du filtre $F$} : convolution de l'entrée $X$
          par le gradient de sortie $\delta O$.
    \item \textbf{Gradient de l'entrée $X$} : corrélation croisée de $\delta O$
          avec le filtre \textbf{retourné de 180°}.
\end{itemize}
Cette symétrie n'est pas un accident : elle découle directement de la structure
de l'opération de convolution. Elle est exploitée par les bibliothèques de
Deep Learning modernes (\textbf{PyTorch}, \textbf{TensorFlow}) pour optimiser
les calculs sur GPU via des opérations \textsc{cuDNN} hautement parallélisées.
\end{resultbox}

\end{document}